\chapter{Notation}

\renewcommand{\arraystretch}{1.3}
\begin{longtable}{|r|l|}
    \hline
        \textbf{Name} & \textbf{Used Symbol} \\ \hline 
        \hline
         Natural Numbers & $\N = \{ 0,1,2, ... \}$ \\ \hline
         Integers & $\Z$ \\ \hline
         Rational Numbers & $\Q$ \\ \hline
         Real Numbers & $\R$ \\ \hline 
         Complex Numbers & $\C$ \\ \hline
         Imaginary Unit & $\i$ \\ \hline
         Complex conjugate & $z^*$ \\ \hline
         Real/Imaginary Part & $\Re z$/ $\Im z$ \\ \hline
         Cartesian product & $\M \times \mathcal{N}$ \\ \hline
         Vectors & $\vec v$ \\ \hline
         Basis vectors & $\vec e_x, \vec e_y, \vec e_z$ \\ \hline
         Matrix & $\mat{A}$ \\ \hline
         Transpose & $\mat{A}^T$ \\ \hline
         Hermitian conjugate & $\mat{A}^\dagger$ \\ \hline
         Diagonal Matrix & $\Diag(v_1, ... v_n) = \begin{pmatrix}
             v_1 \\
             & \ddots \\
             & & v_n
         \end{pmatrix}$ \\ \hline
         Scalar field & $f = f(\vec v)$ \\ \hline
         Vector field & $\vec F = \vec F(\vec v)$ \\ \hline
         Gradient & $\grad f$ \\ \hline
         Gradient w.r.t the variables $\vec v = [x,y,z]$ & $\grad_{\vec v} f$ \\ \hline
         Gradient w.r.t the position of the $n$-th particle & $\grad_n f$ \\ \hline
         Divergence & $\div \vec F$ \\ \hline
         Curl & $\curl \vec F$ \\ \hline
         Hessian & $\Hess f$ \\ \hline
    \caption{Table of Symbols}
    \label{tab:apdx:not:symb}
\end{longtable}
\renewcommand{\arraystretch}{1.0}
When applying a special set of variables,
    the operator will be annotated by a subscript.

\chapter{Single Particle}
\section{Transformation of the Hessian}
\label{apdx:sp:hessian:1d}

We have the Hamiltonian $H(\vec r, \vec p, a_c, a_c^*)$, 
    we cannot calculate the Hessian directly with complex variables,
    since we can only take derivatives with respect to real variables.
To get the Hessian,
    we reformulate the Hamiltonian to $H(\vec r, \vec p, \Re a, \Im a)$,
    which we then can take the Hessian of.
In this formulation we define $\vec\eta' = [\vec r, \vec p, \Re a, \Im a]^T$.
So the Hessian can be formulated as
\begin{equation}
\label{apdx:hess:1d:deletas}
    \delta^2 H = \delta \vec\eta'^\dagger \cdot \Hess_{\vec\eta'}{H} \cdot \delta \vec\eta'.
\end{equation}
We then define a linear transform $\mat{T}$,
    to satisfy $\vec\eta = \mat{T}\cdot\vec\eta'$.
%We then decompose $\mat{T}$ into a transform of $(\Re a, \Im a)\mapsto (a, a^*)$
%    and a permutation that changes the order of the variables.
%We write this as $\mat{T'} = \mat{P}\cdot\mat{T}$.

This fixes the transform $\mat{T}$ to
\begin{equation}
    \mat{T} = \begin{pmatrix}
        \Id_{dim\vec r + \dim\vec p} \\
        & \tfrac12 & - \tfrac{\i}{2} \\
        & \tfrac{1}{2} & \tfrac{\i}{2}
    \end{pmatrix}.
\end{equation}
This transform restores the Wirtinger derivative
\begin{equation}
    \frac{\partial}{\partial z} = 
        \frac{1}{2} 
            \left( \frac{\partial}{\partial\Re z} - \i \frac{\partial}{\partial\Im z}\right).
\end{equation}

%And we can write the permutation as
%\begin{equation}
%    \mat{P} = \begin{pmatrix}
%        \Id_{\dim \vec r}\\
%        & & 1 \\
%        & \Id_{\dim \vec p} \\
%        & & & 1
%    \end{pmatrix}
%\end{equation}

With those definitions we can rewrite the Equation \ref{apdx:hess:1d:deletas} to
\begin{align}
    \delta^2 H & = 
        \delta \vec\eta^\dagger \cdot \left(\mat{T}^{-1} \right)^\dagger
        \cdot \Hess_{\vec\eta'}{H} \cdot 
        \mat{T}^{-1} \cdot\delta \vec\eta,
\end{align}
which gives us the full Hessian for further analysis.

\section{Keeping the Field at its equilibrium}
\label{apdx:sp:fcf}
To check the stability we view the System in a perturbed state $\vec\eta=\vec\eta_0 + \delta\vec\eta$.
We do this for a small perturbation in $z$.
Here we assume that $a$ reaches its equilibrium much faster than $z$,
    such that $a = \alpha_c^{eq}(\vec r)$ always holds.
From Hamilton's equations we immediately get $\dot z = \partial_{p_z} H = p_z/m = 0$.
With this, we know our dynamics are governed by $\ddot z = \dot p_z/m$,
    which we can get from the results in Equation \eqref{eq:sp:z-stabil}.
\begin{equation}
    \ddot z = \frac\alpha m |E_t|^2 L_z.
\end{equation}
If $\ddot z$ points toward the equilibrium, we have a strong indicator that the System is stable and vice versa.
A result of checking this stability can be seen in Figure \ref{fig:apdx:zfasta:stability}.
For this simulation, we assume that the system is stable 
    if the restoring force points opposite to the perturbation.

Here we get the optical bistability we would expect from intuition.
However,
    the approximation is not physical,
    since the cavity needs a finite non-zero relaxation time,
    that will not be instantaneous.
    

\begin{figure}%[h]
    \centering
    %\input{chapters/10_single_cavity/stabilities}
    \includegraphics[width=.8\textwidth]{chapters/99_appendix/stabilities.pdf}
    \caption{
    Plot of the equilibrium position at $\kappa = 180.0 ~ 2\pi kHz$\\ 
    showing the stability of the equilibrium points\\ 
    with the parameters of Table \ref{tab:sp:sim:params}
    }
    \label{fig:apdx:zfasta:stability}
\end{figure}
